\documentclass[11pt,a4paper]{article}
\usepackage{fontspec,amsmath,amssymb,amsthm,microtype}
\setmainfont{lmroman10-regular.otf}[BoldFont=lmroman10-bold.otf,ItalicFont=lmroman10-italic.otf,BoldItalicFont=lmroman10-bolditalic.otf]
\usepackage{xcolor}
\usepackage[most]{tcolorbox}
\usepackage{mathrsfs}
\providecommand{\lt}{<}
\providecommand{\gt}{>}
\setmonofont{DejaVuSansMono.ttf}[BoldFont=DejaVuSansMono-Bold.ttf,ItalicFont=DejaVuSansMono-Oblique.ttf,BoldItalicFont=DejaVuSansMono-BoldOblique.ttf,Scale=0.86]
\newfontfamily\leanfallback{STIXTwoMath-Regular.otf}[Scale=0.86]
\definecolor{leanrule}{gray}{0.62}
\definecolor{leanback}{gray}{0.962}
\definecolor{leantitle}{gray}{0.22}
\definecolor{leanmuted}{gray}{0.38}
\newtcolorbox{leanbox}[1]{enhanced,breakable,colback=leanback,colframe=leanrule,
  boxrule=0.4pt,arc=0.6mm,left=2.4mm,right=2.4mm,top=1.8mm,bottom=1.8mm,
  before skip=0.8em,after skip=1em,fonttitle=\small\bfseries,coltitle=leantitle,
  colbacktitle=leanback,titlerule=0.3pt,title={#1}}
\newlength{\leancharwidth}
\newenvironment{leancode}{\par\smallskip\begingroup\ttfamily\footnotesize\raggedright
  \settowidth{\leancharwidth}{M}\setlength{\parskip}{0pt}\setlength{\parindent}{0pt}
  \frenchspacing}{\par\endgroup\smallskip}
\newcommand{\leanline}[2]{\hangindent=\dimexpr#1\leancharwidth+4\leancharwidth\relax
  \hangafter=1\noindent\hspace*{#1\leancharwidth}#2\par}
\newcommand{\leansource}[1]{\par{\footnotesize\color{leanmuted}#1}\par}
\newcommand{\leanlabel}[1]{\par\smallskip\noindent{\small\bfseries #1}\par}
\usepackage[margin=25mm]{geometry}
\usepackage[unicode,hidelinks]{hyperref}
\hypersetup{pdfauthor={LangevinLean Project}}
\setlength{\emergencystretch}{3em}
\setlength{\parskip}{0.4em}
\setlength{\parindent}{0pt}
\allowdisplaybreaks[2]
\raggedbottom
\title{\texorpdfstring{Gaussian invariance of the velocity thermostat}{Gaussian invariance of the velocity thermostat}}
\author{LangevinLean Project}
\date{}
\begin{document}\maketitle
Draft mathematical exposition; these notes may contain mistakes.


The \(O\) substep removes velocity variance through damping and restores exactly that variance through independent Gaussian noise. This section proves an equality of probability measures, not only a matching second moment. Invariance of this substep is one component of kinetic stationarity; the Hamiltonian evolution and the full diffusion require separate arguments.



The invariant Gaussian temperature is determined by balancing the variance lost through damping against the variance supplied by fresh independent noise. This exact identity is one of the invariant building blocks for kinetic stationarity.


\section{\texorpdfstring{Independent Gaussian linear combinations}{Independent Gaussian linear combinations}}


Let \(G,G\prime\) be independent standard Gaussian vectors in a finite-dimensional real Hilbert space, and suppose \(a^2+b^2=1\). For a continuous linear functional \(L\), the characteristic function of \(aG+bG\prime\) is

\[\exp(-a^2\|L\|^2/2)\exp(-b^2\|L\|^2/2)=\exp(-\|L\|^2/2).\]

Independence is represented by the product of two standard Gaussian measures. The argument constructs the continuous linear map \(A(v,w)=av+bw\), uses the characteristic function of a linear pushforward to pull the test functional back through \(A\), and uses factorization of characteristic functions on independent products to factor its characteristic function. The restrictions of \(L\circ A\) to the two factors are \(aL\) and \(bL\), as follows by evaluation at arbitrary vectors in each factor. The norm of a scalar multiple supplies the squared coefficients. The relation \(a^2+b^2=1\) then simplifies the exponent, and uniqueness of characteristic functions proves the distributional identity.



For independent centered Gaussian vectors \(V\) and \(G\), the characteristic function of \(aV+bG\) factors into the product of the two Gaussian characteristic functions. Its covariance is \((a^2+b^2)I\). Independence is essential: an unspecified cross covariance would contribute an additional term.

\begin{leanbox}{Lean: \texttt{stdGaussian\_\allowbreak{}linear\_\allowbreak{}combination}}
\leansource{\texttt{stdGaussian\_\allowbreak{}linear\_\allowbreak{}combination} \quad KineticLangevin/\allowbreak{}Thermostat.lean, lines 20–42}
\begin{leancode}
\leanline{0}{theorem\ stdGaussian\_linear\_combination\ (a\ b\ :\ ℝ)\ (hab\ :\ a\ \textasciicircum{}\ 2\ +\ b\ \textasciicircum{}\ 2\ =\ 1)\ :}
\leanline{4}{((stdGaussian\ E).prod\ (stdGaussian\ E)).map\ (fun\ z\ ↦\ a\ •\ z.1\ +\ b\ •\ z.2)\ =}
\leanline{6}{stdGaussian\ E}
\end{leancode}
\leanlabel{How the Lean proof follows the argument}
\texttt{std\allowbreak{}Gaussian\_\allowbreak{}linear\_\allowbreak{}combination} gives the formal counterpart of this step. The main intermediate facts are \texttt{A}, \texttt{hleft}, \texttt{hright}, \texttt{heq} in that order. Named proof applications include \texttt{Measure.\allowbreak{}ext\_\allowbreak{}of\_\allowbreak{}char\allowbreak{}Fun\allowbreak{}Dual}, \texttt{heq}; the adjacent proof shows their hypotheses and order of use.
\end{leanbox}

\section{\texorpdfstring{Apply the fluctuation-dissipation balance}{Apply the fluctuation-dissipation balance}}


For nonnegative friction and time step, put \(\delta=e^{-\gamma h/2}\) and \(s=\sqrt{1-\delta^2}\). The variance under the square root is nonnegative, so \(s^2=1-\delta^2\). Hence \(\delta^2+s^2=1\), and the preceding theorem gives \(\delta G+sG\prime\sim\mathcal N(0,I)\). The argument applies the general linear-combination theorem, rewrites the square root squared using its proved nonnegativity, unfolds the variance, and finishes with a polynomial identity. This is precisely the standard-temperature \(O\) half-step used in OBABO.



Take \(a=e^{-\gamma t}\) and \(b=\sqrt{1-e^{-2\gamma t}}\). Their squared sum is one, so a standard Gaussian velocity remains standard Gaussian after thermostat. Position is unchanged, hence the full Gaussian-times-position target is preserved whenever the refresh noise is independent.

\begin{leanbox}{Lean: \texttt{obabo\_\allowbreak{}thermostat\_\allowbreak{}preserves\_\allowbreak{}standardGaussian}}
\leansource{\texttt{obabo\_\allowbreak{}thermostat\_\allowbreak{}preserves\_\allowbreak{}standardGaussian} \quad KineticLangevin/\allowbreak{}Thermostat.lean, lines 46–54}
\begin{leancode}
\leanline{0}{theorem\ obabo\_thermostat\_preserves\_standardGaussian\ \{γ\ h\ :\ ℝ\}}
\leanline{4}{(hγ\ :\ 0\ ≤\ γ)\ (hh\ :\ 0\ ≤\ h)\ :}
\leanline{4}{((stdGaussian\ E).prod\ (stdGaussian\ E)).map}
\leanline{6}{(fun\ z\ ↦\ obaboDamping\ γ\ h\ •\ z.1\ +\ Real.sqrt\ (obaboVariance\ γ\ h)\ •\ z.2)\ =}
\leanline{6}{stdGaussian\ E}
\end{leancode}
\leanlabel{How the Lean proof follows the argument}
\texttt{obabo\_\allowbreak{}thermostat\_\allowbreak{}preserves\_\allowbreak{}standard\allowbreak{}Gaussian} gives the formal counterpart of this step. Named proof applications include \texttt{std\allowbreak{}Gaussian\_\allowbreak{}linear\_\allowbreak{}combination}; the adjacent proof shows their hypotheses and order of use.
\end{leanbox}

\end{document}
